\section{P04 - Évaluation courte - Types construits}

\subsection{Objectifs spécifiques}

\begin{itemize}
\tightlist
\item
  Objectif O1 - Identifier précisément la représentation ou la structure
  en jeu.
\item
  Objectif O2 - Appliquer une méthode disciplinaire complète.
\item
  Objectif O3 - Justifier le résultat sur un cas différent.
\item
  Objectif O4 - Contrôler un cas limite et corriger une erreur observée.
\end{itemize}

\subsection{Capacités officielles atomiques}

\begin{itemize}
\tightlist
\item
  P-DATA-CONSTR-02A
\item
  P-DATA-CONSTR-02D
\end{itemize}

\subsection{Prérequis}

\begin{itemize}
\tightlist
\item
  Reconnaître une consigne liée à tuple.
\item
  Distinguer donnée, méthode et conclusion dans le thème Tuples, listes,
  dictionnaires.
\item
  Rédiger une justification courte en utilisant le vocabulaire du
  programme.
\item
  Contrôler une réponse par un cas limite ou un contre-exemple
  explicite.
\end{itemize}

\subsection{Séance(s) correspondante(s)}

\begin{itemize}
\tightlist
\item
  P04-S1 à P04-S7 : support rattaché aux séances prêtes de la
  progression.
\end{itemize}

\subsection{Situation-problème concrète}

Une station météo stocke des coordonnées fixes, des relevés horaires
modifiables et des mesures accessibles par nom.

\subsection{Activité d’entrée}

\begin{enumerate}
\def\labelenumi{\arabic{enumi}.}
\tightlist
\item
  Identifier ce qui doit rester immuable dans un tuple.
\item
  Modifier une liste de températures.
\item
  Lire une clé dans un dictionnaire de station.
\item
  Décrire ce qui se passe avec une liste vide.
\end{enumerate}

\subsection{Questions}

\subsubsection{Question 1}

\begin{itemize}
\tightlist
\item
  Objectif évalué : O1.
\item
  Capacité officielle : P-DATA-CONSTR-02A.
\item
  Énoncé : résoudre tuple de coordonnées avec \texttt{(36.8,\ 10.2)}.
\item
  Réponse attendue : coordonnées conservées.
\item
  Critère de réussite : méthode visible, résultat correct et contrôle «
  tentative de modification interdite ». \#\#\# Question 2
\item
  Objectif évalué : O2.
\item
  Capacité officielle : P-DATA-CONSTR-02A, P-DATA-CONSTR-02D.
\item
  Énoncé : on dispose de \texttt{releves\ =\ {[}18,\ 20,\ 19{]}}. Écrire
  une boucle \texttt{for\ valeur\ in\ releves} qui calcule la moyenne
  des relevés. Donner la trace (élément courant et cumul à chaque tour).
\item
  Réponse attendue : \texttt{19} (moyenne = (18 + 20 + 19) / 3).
\item
  Critère de réussite : boucle \texttt{for} sur les valeurs (pas les
  indices), résultat correct et contrôle « liste vide ». \#\#\# Question
  3
\item
  Objectif évalué : O3.
\item
  Capacité officielle : P-DATA-CONSTR-02A.
\item
  Énoncé : comparer dictionnaire avec
  \texttt{\{"temperature":\ 21,\ "vent":\ 12\}}.
\item
  Réponse attendue : \texttt{21} pour \texttt{temperature}.
\item
  Critère de réussite : méthode visible, résultat correct et contrôle «
  clé absente ». \#\#\# Question 4
\item
  Objectif évalué : O4.
\item
  Capacité officielle : P-DATA-CONSTR-02A.
\item
  Énoncé : corriger copie de liste pour
  \texttt{{[}{[}1{]},\ {[}2{]}{]}}.
\item
  Réponse attendue : modification locale contrôlée.
\item
  Critère de réussite : méthode visible, résultat correct et contrôle «
  liste imbriquée ». \#\# Barème
\item
  Question 1 : 2 points méthode, 1 point résultat, 1 point justification
  liée à tentative de modification interdite.
\item
  Question 2 : 2 points méthode, 1 point résultat, 1 point justification
  liée à liste vide.
\item
  Question 3 : 2 points méthode, 1 point résultat, 1 point justification
  liée à clé absente.
\item
  Question 4 : 2 points méthode, 1 point résultat, 1 point justification
  liée à liste imbriquée.
\item
  Question 5 : 1 point rôle du tuple, 1 point rôle du dictionnaire, 1
  point rôle de la liste, 1 point calcul du milieu \texttt{(3.0,\ 7.0)}.
  \#\# Erreurs fréquentes
\item
  Erreur fréquente EF1 - Modifier un tuple comme une liste.
\item
  Erreur fréquente EF2 - Parcourir les indices quand les valeurs
  suffisent.
\item
  Erreur fréquente EF3 - Accéder à une clé sans vérifier sa présence.
\item
  Erreur fréquente EF4 - Copier une liste imbriquée seulement au premier
  niveau.
\end{itemize}

\subsection{Remédiation ciblée}

\begin{itemize}
\tightlist
\item
  Activité corrective EF1 : Identifier mutabilité et usage avant
  d'écrire une affectation.
\item
  Activité corrective EF2 : Écrire deux boucles, avec indices puis avec
  valeurs, et comparer.
\item
  Activité corrective EF3 : Tester \texttt{cle\ in\ dictionnaire} avant
  la lecture.
\item
  Activité corrective EF4 : Modifier une sous-liste et observer l'effet
  sur la copie.
\end{itemize}

\subsection{Différenciation}

\begin{itemize}
\tightlist
\item
  Socle : traiter \texttt{(36.8,\ 10.2)} avec une fiche méthode fournie.
\item
  Standard : traiter \texttt{{[}18,\ 20,\ 19{]}} en rédigeant la
  justification complète.
\item
  Expert : inventer un cas limite lié à « clé absente » et expliquer le
  comportement attendu.
\end{itemize}

\subsection{Critères de réussite}

\begin{itemize}
\tightlist
\item
  La capacité officielle est citée dans la copie.
\item
  La méthode contient au moins une étape vérifiable par un pair.
\item
  Le cas limite est discuté avec une donnée concrète.
\item
  La correction explique quelle erreur fréquente est évitée.
\end{itemize}

\subsection{Question complémentaire - structures combinées}

\subsubsection{Question 5}

\begin{itemize}
\tightlist
\item
  Données :
  \texttt{stations\ =\ {[}\{"nom":\ "A",\ "coord":\ (0,\ 4),\ "temperature":\ 21\},\ \{"nom":\ "B",\ "coord":\ (6,\ 10),\ "temperature":\ 18\}{]}}.
\item
  Travail : expliquer le rôle du tuple, de la liste et du dictionnaire
  dans cette donnée, puis calculer le milieu des deux coordonnées.
\end{itemize}

\subsubsection{Corrigé question 5}

\begin{itemize}
\tightlist
\item
  Tuple : chaque \texttt{coord} est un couple immuable.
\item
  Dictionnaire : chaque station associe les clés \texttt{nom},
  \texttt{coord}, \texttt{temperature} à des valeurs.
\item
  Liste : \texttt{stations} rassemble plusieurs dictionnaires et se
  parcourt.
\item
  Résultat attendu : le milieu est
  \texttt{((0+6)/2,\ (4+10)/2)\ =\ (3.0,\ 7.0)}.
\item
  Contrôle : le résultat est un tuple numérique de deux coordonnées.
\end{itemize}

\subsection{Exercices numérotés}

\begin{itemize}
\tightlist
\item
  Exercice 1 : reprendre question 1 en explicitant donnée, méthode,
  résultat et contrôle pour P04.
\item
  Exercice 2 : reprendre question 2 en explicitant donnée, méthode,
  résultat et contrôle pour P04.
\item
  Exercice 3 : reprendre question 3 en explicitant donnée, méthode,
  résultat et contrôle pour P04.
\item
  Exercice 4 : reprendre question 4 en explicitant donnée, méthode,
  résultat et contrôle pour P04.
\end{itemize}

\subsection{Corrigé}

\subsubsection{Corrigé question 1}

\begin{itemize}
\tightlist
\item
  Résultat attendu : coordonnées conservées.
\item
  Méthode exigée : reprendre la démarche du cours puis vérifier le cas
  limite de la question 1.
\item
  Critère de validation : méthode visible, résultat correct et contrôle
  « tentative de modification interdite ». \#\#\# Corrigé question 2
\item
  Résultat attendu : \texttt{19}.
\item
  Méthode exigée : reprendre la démarche du cours puis vérifier le cas
  limite de la question 2.
\item
  Critère de validation : méthode visible, résultat correct et contrôle
  « liste vide ». \#\#\# Corrigé question 3
\item
  Résultat attendu : \texttt{21} pour \texttt{temperature}.
\item
  Méthode exigée : reprendre la démarche du cours puis vérifier le cas
  limite de la question 3.
\item
  Critère de validation : méthode visible, résultat correct et contrôle
  « clé absente ». \#\#\# Corrigé question 4
\item
  Résultat attendu : modification locale contrôlée.
\item
  Méthode exigée : reprendre la démarche du cours puis vérifier le cas
  limite de la question 4.
\item
  Critère de validation : méthode visible, résultat correct et contrôle
  « liste imbriquée ».
\end{itemize}

\subsection{Modalités de passation}

\begin{itemize}
\tightlist
\item
  Durée : 25 minutes.
\item
  Matériel autorisé : fiche types construits, sans corrigé distribué ni
  navigation externe.
\item
  Capacités évaluées :
\item
  P-DATA-CONSTR-02A
\item
  P-DATA-CONSTR-02D
\end{itemize}

\subsection{Fiche liée et aménagement}

\begin{itemize}
\tightlist
\item
  Fiche liée : fiche de cours opérationnelle de la séquence P04, statut
  \texttt{needs\_review}.
\item
  Séance liée : \texttt{P04-S1}, avec question centrée sur tuples,
  listes et dictionnaires.
\item
  Version aménagée : données tuples, listes et dictionnaires surlignées
  et tableau réponse en trois zones.
\item
  Remédiation : corriger une mutation de liste sur un exemple à trois
  éléments, puis verbaliser la méthode en binôme.
\end{itemize}
